% This is text associated with the open textbook "Open Electromagnetics"
% See immediately below for author, date
% This work is licensed under the Creative Commons Attribution-ShareAlike 4.0 International License. (CC BY SA 4.0) 
% To view a copy of the license, visit http://creativecommons.org/licenses/by-sa/4.0/

%ORB TITLE The Maxwell-Faraday Equation 
%ORB AUTHOR 0        # S.W. Ellingson, Virginia Tech (ellingson@vt.edu)
%ORB DATE 20180807   # yyyymmdd
%ORB ABSTRACT The Maxwell-Faraday Equation 

\label{m0050_Maxwell-Faraday_Equation} % <-- Must match name of this file and module directory name.  Don't mess with this!
\index{Maxwell-Faraday equation|(}

In this section, we generalize Kirchoff's Voltage Law (KVL), 
\index{Kirchoff's voltage law!electrostatics}
previously encountered as a principle of electrostatics in Sections~\ref{m0016_Kirchoffs_Voltage_Law_Electrostatics_Integral_Form} and \ref{m0152_KVL_Electrostatics_Integral-to-Differential}.
KVL states that in the absence of a time-varying magnetic flux, the electric potential accumulated by traversing a closed path $\mathcal{C}$ is zero.
Here is that idea in mathematical form:
%
\begin{equation}   
V = \oint_{\mathcal{C}}{ {\bf E} \cdot d{\bf l} } = 0
\end{equation}
%
Now recall Faraday's Law 
\index{Faraday's law}
(Section~\ref{m0055_Faradays_Law}): 
%
\begin{equation}
V = - \frac{\partial}{\partial t} \Phi = - \frac{\partial}{\partial t} \int_{\mathcal{S}} { {\bf B} \cdot d{\bf s} }
\end{equation}
%
Here, $\mathcal{S}$ is any open surface that intersects all magnetic field lines passing through $\mathcal{C}$, with the relative orientations of $\mathcal{C}$ and $d{\bf s}$ determined in the usual way by the Stokes' Theorem 
\index{Stokes' theorem}
convention (Section~\ref{m0051_Stokes_Theorem}).
Note that Faraday's Law agrees with KVL in the magnetostatic case. If magnetic flux is constant, then Faraday's Law says $V=0$. 
However, Faraday's Law is very clearly \emph{not} consistent with KVL if magnetic flux is time-varying.
The correction is simple enough; we can simply set these expressions to be equal.
Here we go:
%
\begin{equation} 
\boxed{  
\oint_{\mathcal{C}}{ {\bf E} \cdot d{\bf l} } = - \frac{\partial}{\partial t} \int_{\mathcal{S}} { {\bf B} \cdot d{\bf s} }
}
\label{m0050_eMFEI}
\end{equation}
%
This general form is known by a variety of names; here we refer to it as the \emph{Maxwell-Faraday Equation} (MFE).

%ORB HIGHLIGHT
\begin{mdframed}[backgroundcolor=yellow!20,nobreak=true] 
The integral form of the Maxwell-Faraday Equation (Equation~\ref{m0050_eMFEI}) states that the electric potential associated with a closed path $\mathcal{C}$ is due entirely to  electromagnetic induction, via Faraday's Law. 
\end{mdframed}

Despite the great significance of this expression as one of Maxwell's Equations, one might argue that all we have done is simply to write Faraday's Law in a slightly more verbose way.  This is true.  The \emph{real} power of the MFE is unleashed when it is expressed in differential, as opposed to integral form.  Let us now do this. 

We can transform the left-hand side of Equation~\ref{m0050_eMFEI} into a integral over $\mathcal{S}$ using Stokes' Theorem.
Applying Stokes' theorem on the left, we obtain 
%
\begin{equation}   
\int_{\mathcal{S}} { \left( \nabla \times {\bf E} \right) \cdot d{\bf s}  }  = - \frac{\partial}{\partial t} \int_{\mathcal{S}} { {\bf B} \cdot d{\bf s} }
\end{equation}
%
Now exchanging the order of integration and differentiation on the right hand side:
%
\begin{equation}   
\int_{\mathcal{S}} { \left( \nabla \times {\bf E} \right) \cdot d{\bf s}  } = \int_{\mathcal{S}} { \left( - \frac{\partial}{\partial t}{\bf B} \right) \cdot d{\bf s} }
\end{equation}
% 
The surface $\mathcal{S}$ on both sides is the same, and we have not constrained $\mathcal{S}$ in any way. $\mathcal{S}$ can be any mathematically-valid open surface anywhere in space, having any size and any orientation.
The only way the above expression can be universally true under these conditions is if the integrands on each side are equal at every point in space.
Therefore,
%
\begin{equation}   
\boxed{
\nabla \times {\bf E} = - \frac{\partial}{\partial t}{\bf B}
}
\label{m0050_eMFED}
\end{equation}
%
which is the MFE in differential form.

What does this mean? Recall that the curl of ${\bf E}$ is a way to take a directive of ${\bf E}$ with respect to position (Section~\ref{m0048_Curl}).  Therefore the MFE constrains spatial derivatives of ${\bf E}$ to be simply related to the rate of change of ${\bf B}$.  Said plainly:

%ORB HIGHLIGHT
\begin{mdframed}[backgroundcolor=yellow!20,nobreak=true] 
The differential form of the Maxwell-Faraday Equation (Equation~\ref{m0050_eMFED}) relates the change in the electric field with position to the change in the magnetic field with time.
\end{mdframed}

Now that \emph{is} arguably new and useful information.  We now see that electric and magnetic fields are coupled not only for line integrals and fluxes, but also at each point in space.  

%%%%%%%%%%%%%%%%%%%%%%%%%%%%%%%%%%%%%%

{\bf Additional Reading:}
\begin{itemize}
\item ``\href{https://en.wikipedia.org/wiki/Faraday's_law_of_induction}{Faraday's Law of Induction}'' on Wikipedia.
\item ``\href{https://en.wikipedia.org/wiki/Maxwell's_equations}{Maxwell's Equations}'' on Wikipedia.
\end{itemize}

\index{Maxwell-Faraday equation|)}


